\section{Normal forms for structural trace}

The previous section ended at the point where every admissible declaration has been replaced
by a canonical normalized public signature.  We now forget the surface presentation and look
only at the structural trace fields that survive this passage to normal form.  This section is
the bridge between the accounting calculus of Sections~2--3 and the cubical depth theorem of
Section~5.

The bridge has one central purpose.  It explains why the later phrase ``higher trace is
derived'' is not merely a syntactic tag.  A normalized trace field is not classified by its
textual shape alone.  It is classified by the semantic role it plays in the sealed interface,
by the historical support on which its typing derivation genuinely depends, and by whether
there is a uniform construction of the field from already exported lower data.

After normalization, every structural trace field generated by the sealing discipline has one
of three roles:
\[
  \text{unary action trace},
  \qquad
  \text{binary comparison trace},
  \qquad
  \text{higher structural horn package}.
\]
The first two roles may contribute irreducible primitive public trace.  The third role may
remain present in the exact realized obligation object, but in the cubical calculus it is
carried by an explicit horn-computation witness and therefore contributes no primitive field
to a \(\mu\)-minimal public signature.

This is the section-level slogan:
\begin{quote}
Normal form does not erase higher structural obligations; it exposes them as horn packages
whose derivedness must be witnessed.
\end{quote}

\subsection{The three trace roles}

Let \(L_{n+1}\) be a sealed layer over a normalized prior state \(B_n\), and let
\(N_{\mathrm{tr}}(L_{n+1})\) be the trace part of its canonical normalized public signature.
A field of \(N_{\mathrm{tr}}(L_{n+1})\) is structural only when its job is to integrate the new
payload with previously exported opaque interface data.  The same mathematical expression
may have a different status when it is authored as new payload.  For example, a path
constructor contributed by a new higher inductive type is payload; a path witnessing that a
new globally acting operation respects an old exported equivalence is structural trace.

\begin{definition}[Normal structural trace role]
A normalized structural trace field \(t\in N_{\mathrm{tr}}(L_{n+1})\) has one of the following
roles.

\begin{enumerate}
\item \textbf{Unary action trace.}  The field records how the new payload acts on one old
opaque public site.  Its historical support contains one prior public interface component
after counting normalization.

\item \textbf{Binary comparison trace.}  The field records compatibility between two public
routes already determined by unary action, transport, old comparison data, or adjacent
composites through those data.  Its historical support contains two independent public
interface components, or an adjacent binary composite whose two components remain visible in
normal form.

\item \textbf{Higher structural horn package.}  The field records the data of a marginal
higher coherence problem generated by sealing.  Its visible boundary is assembled from
already exported unary and binary trace, together with degenerate faces, reindexings, and
transports that counting normalization treats as transparent.  Its missing face and filler are
represented as an open-box extension package.
\end{enumerate}
\end{definition}

The roles should not be read as a hierarchy of textual arities.  A field whose type begins
with a long telescope may still be unary if that telescope ranges over parameters inside one
old opaque interface.  Conversely, a short-looking comparison may be binary if its typing
derivation uses two independent old public exports.  The invariant is historical support after
all prior layers have been collapsed to their opaque normalized signatures.

\begin{table}[t]
\centering
\small
\setlength{\tabcolsep}{3pt}
\begin{tabular}{p{0.17\linewidth}p{0.28\linewidth}p{0.22\linewidth}p{0.21\linewidth}}
\hline
Normal role & Semantic content & Historical support & Primitive status \\
\hline
Unary action trace &
Action of the new payload on one exported interface site. &
One prior public field or one primitive class of an old opaque signature. &
Primitive candidate unless transparently generated. \\
\hline
Binary comparison trace &
Compatibility between two already public routes: actions, transports, equivalences, or
comparison witnesses. &
Two prior public fields, or an adjacent composite trace through them. &
Primitive candidate; necessary in general by the univalent swap obstruction. \\
\hline
Higher structural horn package &
Boundary, missing face, and filler data for a higher structural obligation generated by
sealing. &
A horn boundary assembled from unary and binary public trace plus transparent structural
faces. &
Derived in the fixed cubical calculus when filled by \(\mathsf{hcomp}\), \(\mathsf{transp}\),
and \(\mathsf{hfill}\). \\
\hline
\end{tabular}
\caption{Normal roles for structural trace fields.}
\end{table}

The table is an accounting table, not a claim that only one-dimensional and two-dimensional
mathematics exists.  Higher-dimensional payload may be primitive payload.  What is at issue
is whether a higher field is forced only to make a sealed structural action coherent with old
opaque public data.  Such fields are the horn packages analyzed below and in Section~5.

\subsection{Unary action trace}

Unary action trace is the most direct form of public integration debt.  A globally acting
sealed layer advertises that its payload can be applied to the active interface already
exported by the foundation.  For each active basis site that lies in the declared footprint,
the normalized public signature must contain either a primitive action field or a transparent
construction of that action from previously exported data.

\begin{definition}[Unary action normal form]
A structural trace field \(a\) is in unary action normal form when there are an old primitive
basis site \(s\in S_{\mathrm{bas}}(L_i)\) and a new payload component \(X\) such that, after
canonical normalization, \(a\) has the form
\[
  a_s : \mathsf{Act}_X(s),
\]
up to the telescope isomorphisms and transparent transports quotiented by counting normal
form.  Its support is \(\mathrm{Supp}(a_s)=\{L_i\}\).
\end{definition}

The notation \(\mathsf{Act}_X(s)\) is schematic.  In a universe example it may be an action
on codes; in a modality example it may be a decoding law; in an algebraic interface it may
be the image of a generator.  What matters is that only one old opaque public site is being
acted on.

Unary action trace is usually primitive because later layers cannot reconstruct objectwise
action on an opaque site unless that action has been exported.  A field of unary shape is
nevertheless not counted as primitive if it is transparently generated.  For instance, an
action on a record projection may be definitionally induced by the action on the record
itself, and an action on a rebundled presentation may be a canonical telescope isomorphism
of an already exported action.  The primitive part of the unary trace consists only of the
irreducible classes left after these transparent identifications.

\begin{lemma}[Unary support stability]
If two admissible presentations of a sealed layer are presentation-equivalent, then the unary
action fields of their normalized public trace signatures correspond bijectively up to
primitive-class equivalence.  In particular, presentation equivalence cannot turn an
irreducible unary action on one old opaque basis site into a trace field supported on a
different old layer.
\end{lemma}

\begin{proof}
Presentation equivalence is generated by the normalization moves of Section~3: telescope
rebracketing, record--telescope conversion, definitional expansion and contraction,
transparent transport along already exported equalities, and deletion of fields whose values
are uniformly synthesized.  None of these moves changes which old opaque primitive class is
being acted on.  A generated or rebundled action may be removed from the primitive count,
but an irreducible action on a distinct opaque basis site cannot be transported to another
site unless the normalized old interface already identifies those sites.  The primitive-class
quotient records exactly that identification.
\end{proof}

This lemma is one of the reasons the counting invariant \(\mu\) is computed only after
normalization.  Before normalization, a programmer may write several aliases for the same
action or hide an action inside a record field.  After normalization, the unary action role is
read from the old opaque support and is independent of those packaging choices.

\subsection{Binary comparison trace}

Binary comparison trace is the first genuinely cubical obstruction.  It records that two
public routes agree.  The simplest route comparison has the shape
\[
  r_0,r_1 : P
  \qquad\text{and}\qquad
  c : r_0 = r_1,
\]
but the actual endpoints \(r_0\) and \(r_1\) may themselves be composites of action fields,
old transports, decoding equations, univalence paths, or adjacent comparison trace.

\begin{definition}[Binary comparison normal form]
A structural trace field \(c\) is in binary comparison normal form when, after canonical
normalization, it compares two public composites whose irreducible support uses exactly two
old public components.  Schematically,
\[
  c_{s,t} : \mathsf{Cmp}_X(s,t;r_0,r_1),
\]
where \(s\) and \(t\) are old opaque public sites, \(r_0\) and \(r_1\) are the two normalized
routes built from unary action and previously exported trace, and
\(\mathrm{Supp}(c_{s,t})=\{s,t\}\) after quotienting transparent aliases.
\end{definition}

Binary comparison trace is not redundant in univalent cubical settings.  Objectwise action
on the endpoints of an old equivalence does not by itself prove naturality along that
equivalence.  The two-point swap obstruction of Section~5 makes this precise: the old
interface may contain a public path generated by a nontrivial equivalence, and a sealed
operation must say how its action respects that path.  Therefore binary comparison fields
remain primitive candidates in the normal form.

\begin{remark}[Why binary comparison is not merely two unary actions]
A unary field tells us how \(X\) acts on \(s\), and another unary field tells us how \(X\) acts
on \(t\).  A binary field tells us that these actions commute with a public comparison between
\(s\) and \(t\).  In a univalent foundation the comparison may carry real information: an
equivalence can be transported into an equality of types, and the action of \(X\) must be
compatible with that equality.  Thus the binary role records a naturality square, not a pair
of isolated objectwise values.
\end{remark}

Binary fields are nevertheless subject to the same normalization discipline as unary fields.
A comparison obtained by whiskering an already exported comparison with definitional maps
is derived.  A comparison whose two routes become judgmentally equal after normalization is
derived.  A comparison duplicated by a rebundled record field is not counted twice.  But if
the old public interface contains an irreducible comparison and the new sealed action must
respect it, the resulting naturality witness is primitive unless the derivability relation
exhibits a uniform construction.

\begin{lemma}[Binary comparison support stability]
Presentation equivalence preserves the primitive-class support of irreducible binary
comparison trace.  It may change the displayed routes by canonical isomorphism or transport,
but it cannot eliminate a genuine binary comparison unless the comparison is transparently
constructed from already exported lower trace.
\end{lemma}

\begin{proof}
The proof is the same normal-form argument as for unary support, with one extra observation.
The two endpoints of a binary comparison are route expressions in the normalized public
interface.  Presentation moves may reassociate, unfold, or transport these route expressions,
but after quotienting by transparent operations the pair of independent old supports remains
visible.  If the comparison disappears, then the presentation relation must have supplied a
specific transparent construction of it from previously exported data.  In that case the field
is derived by definition of the primitive/derived quotient and contributes no primitive trace.
Otherwise it remains an irreducible binary class.
\end{proof}

\subsection{Higher trace as a structural horn package}

The third normal role is the one used by the upper-bound proof.  A higher structural field is
not represented as an arbitrary choice of a filler for a completely fixed closed boundary.  It
is represented as a marginal extension problem.  The lower faces are already public.  One
remote comparison face is not yet public.  The obligation is to produce that missing face and
a filler making the resulting open box commute.

This distinction is crucial in a cubical setting.  The space of fillers of a closed higher
boundary need not be contractible.  Cubical types may have genuine higher homotopy.  The
paper therefore does not claim that all higher fillers are unique.  It claims that the
sealing-generated higher structural obligation has open-box extension form, and that the
fixed cubical operations compute the missing structural data needed for public normalization.

\begin{definition}[Structural horn package in normal form]
A normalized higher structural trace field \(h\) is a structural horn package when its typing
data can be decomposed as
\[
  h \equiv (\partial_h,\gamma_h,F_h),
\]
where:
\begin{enumerate}
\item \(\partial_h\) is the visible boundary assembled from normalized unary action trace,
normalized binary comparison trace, and transparent degenerate or transport faces;
\item \(\gamma_h\) is the missing remote comparison face compatible with the endpoints of
\(\partial_h\);
\item \(F_h\) is the filler witnessing that \(\partial_h\) together with \(\gamma_h\) forms the
required higher structural cell.
\end{enumerate}
Equivalently, \(h\) presents an element of a horn-extension telescope
\[
  \sum_{\gamma:\Gamma_{\mathrm{miss}}(\partial_h)}
    \mathsf{Fill}_{\partial_h}(\gamma).
\]
\end{definition}

The notation is intentionally close to the later definition of \(H_k(\partial)\).  This
section classifies the normal role; Section~5 proves that the corresponding cubical
open-box extension object is contractible and substitution-stable in the fixed calculus.

\begin{remark}[The boundary is public]
The boundary \(\partial_h\) is not hidden derivation state.  It is built from the normalized
public signature exported by earlier layers and from the lower-arity trace of the current
sealed layer.  This is why a later layer can reuse the horn computation without reopening the
opaque proof that produced the current layer.  If the boundary of a proposed higher field
depends on private construction details not present in the normalized public interface, then
that field is not a well-formed structural horn package for the present calculus.
\end{remark}

A useful mental picture is the first nontrivial case.  Suppose the new layer has unary action
on a recent public site and binary comparison against an adjacent public comparison.  When a
further remote layer is exposed, the already exported unary and binary data determine all
but one face of a three-dimensional compatibility box.  The remaining face is the direct
remote comparison.  The structural horn package asks for that face together with a filler.
The cubical theorem later says that this package is computed from the displayed boundary by
the ambient Kan operations.

\subsection{Derivedness witnesses}

The word ``derived'' is safe only when accompanied by a witness.  In this paper a higher
structural field is derived because there is a uniform computation that reconstructs it from
its lower boundary.  The witness is part of the normalized role, and the proof of its adequacy
is discharged by the open-box theorem of Section~5.

\begin{definition}[Semantic derivedness witness]
Let \(h\) be a structural horn package with visible boundary \(\partial_h\).  A
\emph{semantic derivedness witness} for \(h\) consists of:
\begin{enumerate}
\item a cubical term
\[
  t_h(\partial_h)
    : \sum_{\gamma:\Gamma_{\mathrm{miss}}(\partial_h)}
        \mathsf{Fill}_{\partial_h}(\gamma),
\]
constructed uniformly from the visible boundary using the structural operations admitted by
the fixed calculus, in particular composition, transport, and filling;
\item boundary agreement equations showing that the computed missing face and filler restrict
to the displayed visible faces;
\item substitution stability: for every substitution \(\sigma\), the reindexed computation
\(\sigma^*t_h(\partial_h)\) agrees with the computation on the reindexed boundary
\(t_h(\sigma^*\partial_h)\), up to the canonical equality of the cubical structure;

\item \emph{Public normalization.}  Replacing an explicitly presented field \(h\) by
\(t_h(\partial_h)\) gives a presentation-equivalent normalized public signature in which
\(h\) is not a primitive trace field.
\end{enumerate}
\end{definition}

The term \(t_h\) is the formal content behind the phrase ``computed horn''.  In the fixed
cubical calculus, Section~5 constructs \(t_h\) by turning the structural horn package into an
open-box extension problem and applying the cubical composition and filling operations.  The
transport component aligns the computed face with the already exported comparison data.
The filling component supplies the interior cell.

\begin{remark}[Derived is not a color on syntax]
It would be unsatisfactory to define a raw syntax constructor
\(\mathsf{derivedHorn}\) and then conclude that all higher trace is derived.  The normal-form
judgment used here is stronger.  A higher field may be assigned the horn role only when its
boundary is public and its derivedness obligation is a theorem-facing computation problem.
For the fixed cubical calculus, the computation is discharged by the open-box theorem.  For a
different surface language or foundation, the adequacy package must preserve this witness.
\end{remark}

This is also where the paper's two notions of depth begin to separate.  A horn package may
remain present as an exact obligation factor: the realized object still contains the type of
missing-face-and-filler data.  Once a semantic derivedness witness has been supplied, however,
that package is absent from the primitive part of the normalized public signature.  Thus
higher exact factors can be present and contractible while higher primitive public fields are
eliminated.

\begin{proposition}[Primitive status of the three roles]
Assume every normalized higher structural horn package in a sealed declaration carries a
semantic derivedness witness.  Then the primitive part of the normalized structural trace
signature is contained in the union of unary action trace and binary comparison trace:
\[
  \operatorname{Prim}\bigl(N_{\mathrm{tr}}(L_{n+1})\bigr)
  \subseteq
  \operatorname{Act}_{\mathrm{prim}}(L_{n+1})
  \sqcup
  \operatorname{Cmp}_{\mathrm{prim}}(L_{n+1}).
\]
Equivalently, no higher structural horn package contributes an irreducible primitive public
trace field to \(\mu\).
\end{proposition}

\begin{proof}
By the trichotomy of normal roles, each normalized structural trace field is unary, binary,
or horn-shaped.  Unary and binary fields remain primitive exactly when they are not removed
by transparent generation, canonical telescope isomorphism, or presentation equivalence.  A
horn-shaped field with a semantic derivedness witness is uniformly replaceable by
\(t_h(\partial_h)\).  Public-normalization compatibility makes the presentation with the
explicit field equivalent to the presentation in which the field is synthesized.  Counting
normal form counts only irreducible primitive schemas, so the horn field contributes zero to
the primitive trace cardinality.
\end{proof}

The proposition is deliberately conditional.  Section~4 supplies the normal-form
classification and states what a derivedness witness must contain.  Section~5 proves that the
structural horn packages generated by the fixed cubical sealed-extension calculus have such
witnesses.

\subsection{Payload is not reclassified as trace}

The theorem is about public structural integration obligations, not about a collapse of
higher mathematics.  This distinction is easiest to lose precisely at the point where higher
horn packages appear.  The normal-form discipline therefore includes a payload exclusion
rule.

\begin{definition}[Payload exclusion]
A field belongs to payload rather than structural trace when it is authored as new semantic
content of the sealed layer and is not forced merely by the requirement that the new payload
act coherently on the old opaque interface.  In particular, the following are payload when
introduced as part of the new layer:
\begin{enumerate}
\item new operations, type formers, modalities, or universe constructors;
\item algebraic laws internal to the new operation;
\item higher inductive constructors or higher path constructors contributed as new
mathematical structure;
\item user-specified higher cells whose boundary is internal to the new payload rather than
assembled from prior public trace.
\end{enumerate}
Such fields are counted by the payload invariant \(\kappa\), not by the structural trace cost
\(\mu\), unless the sealing discipline separately requires public integration trace for their
action on old interfaces.
\end{definition}

The exclusion rule prevents a misleading reading of the depth-two theorem.  A layer may
export a primitive three-dimensional constructor as payload.  The theorem says only that the
additional structural coherence required to integrate that constructor with the existing
sealed interface is normalized as unary action, binary comparison, and derived horn
computation.  It does not say that the constructor itself is derived or absent.

\begin{lemma}[No payload-to-trace leakage]
Canonical normalization preserves the distinction between payload fields and structural trace
fields.  A payload field can be accompanied by structural trace explaining its action on old
interfaces, but the payload field itself is not moved into \(N_{\mathrm{tr}}\) merely because it
has higher arity or higher dimension.
\end{lemma}

\begin{proof}
The normalized public signature decomposes into payload and trace parts before primitive
trace cost is computed.  Presentation moves may rebundle a payload field, unfold a record
containing it, or transport it along an exported definitional equality, but they do not change
its semantic role as new content of the sealed layer.  A field enters the trace part only when
its typing judgment is an integration obligation against previously exported opaque data.
Thus higher-dimensional payload may increase \(\kappa\), while the induced public
compatibility fields are classified separately by the unary/binary/horn trichotomy.
\end{proof}

\subsection{Role uniqueness and presentation invariance}

The normal roles are useful only if they are stable under the presentation moves already
quotiented by the adequacy bridge.  The following proposition records the invariant that is
used implicitly whenever later sections reason about ``the'' action part, comparison part, or
horn part of a normalized trace signature.

\begin{proposition}[Trace-role normal form]
Every structural field in the canonical normalized trace signature of an admissible sealed
layer has a unique normal role, up to the primitive-class quotient of Section~2 and the
presentation equivalence of Section~3.  Presentation-equivalent declarations have canonically
isomorphic role decompositions:
\[
  N_{\mathrm{tr}}(e)
  \cong
  N_{\mathrm{tr}}(e')
  \quad\Longrightarrow\quad
  \begin{cases}
    \operatorname{Act}(e) \cong \operatorname{Act}(e'),\\
    \operatorname{Cmp}(e) \cong \operatorname{Cmp}(e'),\\
    \operatorname{Horn}(e) \cong \operatorname{Horn}(e').
  \end{cases}
\]
The isomorphisms preserve historical support and primitive/derived status.
\end{proposition}

\begin{proof}
Canonical normalization first flattens the exported public telescope, then quotients by the
transparent presentation moves, and finally records primitive classes.  The role of a
structural field is determined by the irreducible support of its typing derivation in this
normal form.  Support of size one gives unary action; support of size two and route-comparison
shape gives binary comparison; the inductive sealing grammar gives horn packages when a
higher marginal coherence problem is presented as a public boundary plus missing face and
filler.  Since presentation equivalence preserves support, route endpoints up to canonical
isomorphism, and the primitive/derived tag relation, it also preserves the role decomposition.
Uniqueness follows because the three cases are disjoint after opacity: a field cannot both
act on one old opaque site and compare two independent old routes, and a horn package is
recognized by the presence of a public boundary plus missing-face-and-filler telescope rather
than by a single route equality.
\end{proof}

\begin{remark}[Adjacent composites]
The phrase ``support of size two'' should be read after normalization, not by counting the
number of names in a displayed term.  An adjacent composite may mention several textual
components, but if all but two are transparent degeneracies or transports generated from the
normalized public signature, the field is binary.  Conversely, a term with few variables may
still be higher if its typing derivation depends irreducibly on a horn boundary assembled from
several independent prior exports.
\end{remark}

\subsection{A running example of the roles}

The \(\mathsf{SMod}\) instance of Section~3 illustrates the normal-form classification.  The
modal payload \(\mathsf{Flat}\) and its unit are payload.  The action on old universe codes
and the decoding comparison are unary trace.  The compatibility with function codes and
function types is binary comparison trace.  The route coherence comparing the two ways to
transport the modal action through the old decoding law is a structural horn package.

Schematically, the two routes in the function-code case have the form
\[
  \mathsf{El}(\mathsf{FlatU}(\mathsf{Arr}(A,B)))
  \rightrightarrows
  \mathsf{Flat}(\mathsf{El}(A)) \to \mathsf{Flat}(\mathsf{El}(B)).
\]
The route endpoints are assembled from unary decoding trace and binary function-comparison
trace.  The higher field comparing the routes does not introduce a new primitive interaction
site.  Its boundary is public, and the later open-box theorem supplies the missing
structural computation.

This example also shows why the normal-form roles should be assigned after adequacy rather
than from surface syntax.  A surface declaration may write the route coherence explicitly,
omit it and rely on elaboration, or hide it inside a record of modal laws.  These
presentations differ textually but normalize to the same trace role decomposition: the
coherence is a horn package, not an additional primitive public comparison.

\subsection{How this section is used later}

The depth theorem in Section~5 uses this section in two directions.  For the upper bound,
the normal form isolates all possible higher structural obligations as horn packages.  The
cubical open-box theorem then supplies derivedness witnesses for those packages, so they
contribute no primitive trace above binary arity and their exact realized factors are
contractible.  For the lower bound, the normal form isolates binary comparison trace as the
first place where univalence can create an irreducible obstruction.  The boolean-swap example
therefore targets exactly the binary role and shows that unary trace cannot be sufficient in
general.

The recurrence theorem in Section~6 uses a different consequence.  Once the primitive trace
signature has been reduced to unary and binary roles, a separate recent-history theorem can
ask which layers those binary comparisons must mention.  Arity depth two alone does not give
a chronological window.  The role normal form says only that primitive structural trace has
unary or binary support; the later factorization theorem says that surviving binary support
can be chosen among the two most recent normalized signatures under the fixed export
hypotheses.

Thus Section~4 should be read as the classification theorem for trace.  Section~5 supplies
the cubical computation that makes the horn role derived.  Section~6 counts the resulting
primitive unary and binary public fields under additional coupling and chronology hypotheses.
